% =====================================================================
% §7.5  Grounding mechanism: output-interface granularity + semantic-field
%        necessity + verifier role
% Authoritative claim set: M1, M3, L3. L0 = full STC-Bench frame (n=2000), camera-ready final.
% =====================================================================
\subsection{What makes the interface work: granularity, semantic fields, verifier}
\label{sec:mechanism}

\paragraph{Output-interface granularity (L0/L1/L2).}
We ablate the grounding \emph{output interface} on a common base (Qwen2.5-1.5B), corpus
(\texttt{fixv2\_interface\_align}), budget (matched LoRA-lite), and seed (1337), varying only how the
grounded constraint is emitted:
\begin{itemize}
  \item \textbf{L1 dense} — a weight for every edge ($\Theta(|E|)$). This is Dense-LoRA: it collapses to
        a fixed-capacity, scenario-invariant enumeration that neither scales with $|E|$ nor encodes the
        constraint (CS $0.02$, \S\ref{sec:closedloop}).
  \item \textbf{L2 typed sparse anchors} — STAG: \texttt{road$\mid$dir$\mid$range$\mid$type$\mid$level}
        per constraint (CS $0.96$).
  \item \textbf{L0 untyped road-level} — road$+$range only, with the typed semantic fields removed
        (single varied variable vs.\ L2 = typed structure). On the full STC-Bench frame ($n{=}2000$;
        $1840$ active-constraint cases), L0 is \emph{not} degenerate: on a frame-free road-text metric it
        names the correct roads \emph{as well as} L2 (road-level recall $0.869$ vs.\ L2's $0.868$), with
        schema $1.0$ and $0$ hallucination, but, lacking typed output, collapses on the routing-relevant
        axes (direction $0.16$, conflict resolution $0.00$, vs.\ L2's $0.99$/$0.75$); its typed
        anchor-recall is $15.95$ \emph{by interface design} (vs.\ L2 $95.70$).
\end{itemize}
The three sit monotonically on output structure (dense $\to$ typed-sparse $\to$ untyped). L1 fails to
localize (CS $0.02$); L0 under-structures (finds roads but cannot direct or resolve conflicts); only L2 typed
sparse anchors both stay selective and carry the semantics routing needs (CS $0.96$). The result isolates
\emph{typed structure}, not mere sparsity, as the load-bearing property.
% T8c canonical P2 mechanism (Full; semantics locked per p2-final-mechanism-locked memory)
Mechanistically, the dense interface emits a fixed-capacity ($\sim$130--210-slot), near-uniform,
scenario-invariant edge enumeration: it neither scales with $|E|$---at $|E|{=}398$ it emits $\sim$206 of
$398$ edges before terminating voluntarily, well under the token budget---nor encodes the instructed
constraint, since its weights are near-uniform and invariant across events; changed-edge recall therefore
decays as $|E|$ outgrows this capacity and constraint satisfaction stays $\approx 0$, while STAG's typed
sparse anchors ($k\approx 1$) scale trivially and carry the constraint typing.
% Output-interface granularity ablation (matched LoRA-lite, Qwen2.5-1.5B, seed1337).
% FULL STC-Bench frame (n=2000; 1840 active-constraint cases).
% Road recall = frame-free road-text recall (run_diagnostic convention), computed identically for
% L0 and L2 -> directly comparable. L0/L2 typed anchor-recall (15.95 / 95.70) reported in text.
% Source: aaai27_camera_ready_full_frame/{l0_eval_summary_full.json, eval/summary_qwen1p5b_full.json}.
\begin{table}[t]
\centering
\small
\caption{Output-interface granularity. Same base/corpus/budget/seed; only the grounding output
representation varies. On a frame-free road-text metric L0 and L2 name roads \emph{equally well}
($0.87$); but lacking typed structure L0 cannot express the direction/conflict semantics routing requires;
L1 fails to localize; typed sparse anchors (L2) are the operating point. ($^\dagger$Road recall is frame-free
road-text recall, scorer-frame-independent, computed identically for L0 and L2; L0's typed anchor-recall
is low \emph{by interface design} ($15.95$ vs.\ L2 $95.70$). Full STC-Bench frame, $n{=}2000$.)}
\label{tab:granularity}
\begin{tabular}{lccccc}
\toprule
Interface & Road recall$^\dagger$ & Direction & Conflict res. & Schema & Closed-loop CS \\
\midrule
L1 dense edge-weight   & --    & --   & --   & --   & 0.02 (uniform enum.) \\
L0 untyped road-level & 0.87 & 0.16 & 0.00 & 1.00 & -- \\
\textbf{L2 typed sparse (ours)}  & 0.87 & 0.99 & 0.75 & 1.00 & \textbf{0.96} \\
\bottomrule
\end{tabular}
\end{table}
  % L0/L1/L2 full-frame final (n=2000)

\paragraph{Which semantic fields carry the grounding (P5 field ablation).}
Masking individual input fields (length-preserving mask token; canonical STAG; $n{=}180$) isolates each
field's load:
\begin{itemize}
  \item \textbf{Direction} is the largest lever for \emph{localization}: anchor recall $100\to 51.11$
        ($-48.89$), and conflict resolution $100\to 57.5$.
  \item \textbf{Event-type/severity} is the dominant field for \emph{conflict resolution}: $100\to 50.0$.
  \item Temporal masking costs little ($93.33$ anchor recall), and the explicit conflict directive is
        \emph{redundant} (0 drop): STAG resolves conflicts from event$+$temporal content, not the literal
        ``override'' cue---a robustness finding, not a gap.
\end{itemize}
These dependencies mirror the closed-loop and offline frontiers (direction/conflict are where difficulty
concentrates).

\paragraph{The verifier normalizes, it does not repair.}
Comparing with vs.\ without the verifier (canonical, $n{=}2000$): machine-consumable/schema validity rises
$83.45\to 100$, but this gain is \emph{entirely} from normalizing correct \texttt{NO\_ANCHOR} sentinels
(16.55\% of cases) into the uniform schema---the raw model already emits \textbf{0 malformed JSON and 0
hallucinated roads}. Route/graph validity ($99.4$/$100$) are guarantees the verifier adds (the raw output
has no reachability check); a directional-consistency guard catches $5$ residual direction errors
(leave-one-out, $n{=}640$). The verifier delivers deployability (uniform schema $+$ reachability
guarantees) on top of an already-clean model, rather than repairing frequent model errors.
